% ============ §6 动词 ============
\hdr{§6}{动词：两个变位组，三个时态}

\callout{\textbf{这一节的坐标}：动现 p3--p16（两组变位、-ся 动词）、p26--p35（язы́к 与副词）；动过 p2--p13（-л 型过去时）、p15--p19（时间词与星期）；动将 p2--p8（бу́ду ＋ 不定式）。\\
\textbf{这一节要解决的事}：给一个动词，能说出它属于哪一组、六个现在时人称怎么写、过去时四个形式怎么写、将来时怎么说。}

\subhdr{6.1　两个变位组：先看不定式词尾}

动现 p3 的分类只有两行：\textbf{-ать／-ять} 是第一组，\textbf{-ить／-еть} 是第二组。剩下的差别全在词尾上。表头写的 \textbf{чита́ть}、\textbf{говори́ть} 是\textbf{不定式}，也就是动词的\textbf{原形}——词典、单词表里只收这一个形式；下面六行的形式都由它变出来：

\begin{minipage}{\linewidth}
\tabtop
\begin{tabular}{@{}>{\raggedright\arraybackslash}p{0.12\linewidth}>{\raggedright\arraybackslash}p{0.22\linewidth}>{\raggedright\arraybackslash}p{0.24\linewidth}>{\raggedright\arraybackslash}p{0.32\linewidth}@{}}
\thead{人称} & \thead{第一组 чита́ть} & \thead{第二组 говори́ть} & \thead{特殊：писа́ть（第一组）} \\
\end{tabular}
\tabmid
\begin{tabular}{@{}>{\raggedright\arraybackslash}p{0.12\linewidth}>{\raggedright\arraybackslash}p{0.22\linewidth}>{\raggedright\arraybackslash}p{0.24\linewidth}>{\raggedright\arraybackslash}p{0.32\linewidth}@{}}
я & чита́ю & говорю́ & пишу́ \\
ты & чита́ешь & говори́шь & пи́шешь \\
он／она́ & чита́ет & говори́т & пи́шет \\
мы & чита́ем & говори́м & пи́шем \\
вы & чита́ете & говори́те & пи́шете \\
они́ & чита́ют & говоря́т & пи́шут \\
\end{tabular}
\tabbot
\end{minipage}

\textbf{这张表在说什么}：两组的分界线在\textbf{词尾的第一个元音}——第一组是 е／ю（чита́ешь、чита́ют），第二组是 и／я（говори́шь、говоря́т）。只有一个例外：词干末尾是 \textbf{ж、ш、ч、щ} 时，第一人称单数写 \textbf{-у} 不写 \textbf{-ю}、第三人称复数写 \textbf{-ут} 不写 \textbf{-ют}。

\textbf{逐行读法}：第四列是这一节的第一个坑——\textbf{писа́ть} 词尾是 -ать（第一组），但词干 пис- 后面接的是 ш，所以 \textbf{я пишу́}（不是 *писа́ю）、\textbf{они́ пи́шут}（不是 *писа́ют）。另外两个要单独记：\textbf{спать}（课件在第二组里打了星号）→ \textbf{сплю／спят}；\textbf{гото́вить} → \textbf{гото́влю／гото́вят}（词干末尾 л 之后接 ю）。

\aside{这条「ж／ш／ч／щ 后写 у 不写 ю」与 \S4.3 的「г／к／х／ж／ш／ч／щ 后写 и 不写 ы」是同一条正字法规律的两次出现：这七个字母的后面，元音要「换一个」。所以它不是动词的特例，而是拼写法的通例。}

\textbf{课件点名的「难动词」（动将 p2--p3）}。这七个词是变位里最容易写错的一批，单独列一遍：

\begin{minipage}{\linewidth}
\tabtop
\begin{tabular}{@{}>{\raggedright\arraybackslash}p{0.18\linewidth}>{\raggedright\arraybackslash}p{0.16\linewidth}>{\raggedright\arraybackslash}p{0.16\linewidth}>{\raggedright\arraybackslash}p{0.34\linewidth}@{}}
\thead{不定式} & \thead{я} & \thead{они́} & \thead{难在哪} \\
\end{tabular}
\tabmid
\begin{tabular}{@{}>{\raggedright\arraybackslash}p{0.18\linewidth}>{\raggedright\arraybackslash}p{0.16\linewidth}>{\raggedright\arraybackslash}p{0.16\linewidth}>{\raggedright\arraybackslash}p{0.34\linewidth}@{}}
писа́ть & пишу́ & пи́шут & ш 后写 -у／-ут \\
переводи́ть & перевожу́ & перево́дят & д → ж（я 形） \\
занима́ться & занима́юсь & занима́ются & 元音后 -сь、辅音后 -ся \\
спать & сплю & спят & 词干加 л \\
гото́вить & гото́влю & гото́вят & 词干加 л，-вить 型 \\
пить & пью & пьют & 词干 пь-，不是 *пи́ю \\
есть & ем & едя́т & 变位不规则（ем／ешь／ест／еди́м／еди́те／едя́т） \\
\end{tabular}
\tabbot
\end{minipage}

\textbf{这张表在说什么}：七个词分别踩了四个不同的坑——拼写（писа́ть）、词干交替（переводи́ть 的 д→ж）、反身后缀（занима́ться）、词干增音（спать、гото́вить、пить、есть）。它们不构成新规则，只是把 \S6.1 那条规律落实到具体词上。

\subhdr{6.2　-ся 动词：元音后写 -сь，辅音后写 -ся}

动现 p8 与 p16 给了 \textbf{занима́ться}、\textbf{учи́ться} 两个词。它们的词尾跟普通动词一样，只是后面拖一个小尾巴：

\begin{minipage}{\linewidth}
\tabtop
\begin{tabular}{@{}>{\raggedright\arraybackslash}p{0.12\linewidth}>{\raggedright\arraybackslash}p{0.22\linewidth}>{\raggedright\arraybackslash}p{0.22\linewidth}>{\raggedright\arraybackslash}p{0.30\linewidth}@{}}
\thead{人称} & \thead{занима́ться} & \thead{учи́ться} & \thead{尾巴怎么写} \\
\end{tabular}
\tabmid
\begin{tabular}{@{}>{\raggedright\arraybackslash}p{0.12\linewidth}>{\raggedright\arraybackslash}p{0.22\linewidth}>{\raggedright\arraybackslash}p{0.22\linewidth}>{\raggedright\arraybackslash}p{0.30\linewidth}@{}}
я & занима́юсь & учу́сь & 词尾是元音 → -сь \\
ты & занима́ешься & у́чишься & 词尾是辅音 → -ся \\
он／она́ & занима́ется & у́чится & -ет／-ит ＋ ся \\
мы & занима́емся & у́чимся & -ем／-им ＋ ся \\
вы & занима́етесь & у́читесь & -ете／-ите ＋ сь \\
они́ & занима́ются & у́чатся & -ют／-ят ＋ ся \\
\end{tabular}
\tabbot
\end{minipage}

\textbf{这张表在说什么}：一行一条规则——\textbf{词尾以元音收尾就写 -сь，以辅音收尾就写 -ся}。所以 я／вы 用 -сь（занима́юсь、занима́етесь），其余四个用 -ся。

\textbf{逐行读法}：第二列与第三列不是两个变位组的问题，занима́ться 是第一组（-ать）、учи́ться 是第二组（-ить），所以词尾一个用 -ю／-ешь、一个用 -ю／-ишь；尾巴的规则两列相同。

\textbf{补：读音}。\textbf{-тся} 读 [ца]（у́чится 读 [у́чица]），\textbf{-ться} 也读 [ца]（учи́ться 读 [учи́ца]）。课件在 У8 p12 的 \textbf{нра́вится} 旁标了 [нрав'иц:а]，就是这个现象。

\subhdr{6.3　过去时：不看人称，只看性数}

动过 p2--p4 用一句话说完了整个过去时：\textbf{把不定式的 -ть 去掉，加 -л／-ла／-ло／-ли}。

\begin{minipage}{\linewidth}
\tabtop
\begin{tabular}{@{}>{\raggedright\arraybackslash}p{0.24\linewidth}>{\raggedright\arraybackslash}p{0.14\linewidth}>{\raggedright\arraybackslash}p{0.20\linewidth}>{\raggedright\arraybackslash}p{0.18\linewidth}>{\raggedright\arraybackslash}p{0.14\linewidth}@{}}
\thead{谁} & \thead{词尾} & \thead{чита́ть} & \thead{быть} & \thead{учи́ться} \\
\end{tabular}
\tabmid
\begin{tabular}{@{}>{\raggedright\arraybackslash}p{0.24\linewidth}>{\raggedright\arraybackslash}p{0.14\linewidth}>{\raggedright\arraybackslash}p{0.20\linewidth}>{\raggedright\arraybackslash}p{0.18\linewidth}>{\raggedright\arraybackslash}p{0.14\linewidth}@{}}
он（或男性说话人的 я） & -л & чита́л & был & учи́лся \\
она́（或女性说话人的 я） & -ла & чита́ла & была́ & учи́лась \\
оно́ & -ло & чита́ло & бы́ло & учи́лось \\
мы／вы／они́ & -ли & чита́ли & бы́ли & учи́лись \\
\end{tabular}
\tabbot
\end{minipage}

\textbf{这张表在说什么}：过去时\textbf{不跟人称变，只跟「谁」的性和数变}。同一个 я，男的写 \textbf{чита́л}、女的写 \textbf{чита́ла}——这不是语法错误与否的问题，而是\textbf{说话人自己的性别}决定。

\textbf{逐行读法}：第四列 быть 是最不规则的一个（был／была́／бы́ло／бы́ли，阴性重音跑到词尾）；第五列 учи́ться 带 -ся，过去时用 -лся／-лась／-лось／-лись，阳性那个 -лся 里的 л 后面直接接 ся。

\aside{课件没有解释「为什么不看人称」。原因是俄语过去时来自古代的过去分词，分词只跟性数一致，所以人称信息在历史演变中丢掉了。记住结论就够：\textbf{现在时看人称，过去时看性数}。}

\subhdr{6.4　将来时：бу́ду ＋ 不定式，只有一个公式}

动将 p8 的将来时比现在时简单得多——\textbf{бу́ду 按人称变，动词保持不定式不变}：

\begin{minipage}{\linewidth}
\tabtop
\begin{tabular}{@{}>{\raggedright\arraybackslash}p{0.10\linewidth}>{\raggedright\arraybackslash}p{0.20\linewidth}>{\raggedright\arraybackslash}p{0.56\linewidth}@{}}
\thead{人称} & \thead{бу́ду 的形式} & \thead{例（动将 p8）} \\
\end{tabular}
\tabmid
\begin{tabular}{@{}>{\raggedright\arraybackslash}p{0.10\linewidth}>{\raggedright\arraybackslash}p{0.20\linewidth}>{\raggedright\arraybackslash}p{0.56\linewidth}@{}}
я & бу́ду & За́втра я бу́ду повторя́ть слова́. \\
ты & бу́дешь & Ты бу́дешь чита́ть. \\
он／она́ & бу́дет & Он бу́дет отдыха́ть ле́том. \\
мы & бу́дем & Мы бу́дем у́жинать. \\
вы & бу́дете & Ско́ро вы бу́дете хорошо́ понима́ть по-ру́сски. \\
они́ & бу́дут & Но́чью студе́нты бу́дут спать. \\
\end{tabular}
\tabbot
\end{minipage}

\textbf{这张表在说什么}：六个形式全部来自同一个词根 буд-，加的是第一组词尾（-у／-ешь／-ет／-ем／-ете／-ут）。动词部分永远是 \textbf{-ть} 结尾的不定式——\textbf{不变位}。

\textbf{逐行读法}：第三行 он 与 она́ 用同一个 бу́дет；第五行的例子里 \textbf{понима́ть} 保持原形，如果换成现在时就得写 \textbf{понима́ете}——这一个对比就能看出将来时的构造。

\subhdr{6.5　三时态对照：一句话三个时间}

动将 p6--p8 与 p12 把三个时态排在一起，配上时间副词就成了一张完整的表：

\begin{minipage}{\linewidth}
\tabtop
\begin{tabular}{@{}>{\raggedright\arraybackslash}p{0.14\linewidth}>{\raggedright\arraybackslash}p{0.23\linewidth}>{\raggedright\arraybackslash}p{0.28\linewidth}>{\raggedright\arraybackslash}p{0.23\linewidth}@{}}
\thead{时间} & \thead{怎么构成} & \thead{时间副词} & \thead{例} \\
\end{tabular}
\tabmid
\begin{tabular}{@{}>{\raggedright\arraybackslash}p{0.14\linewidth}>{\raggedright\arraybackslash}p{0.23\linewidth}>{\raggedright\arraybackslash}p{0.28\linewidth}>{\raggedright\arraybackslash}p{0.23\linewidth}@{}}
过去时 & 词干 ＋ -л／-ла／-ло／-ли & вчера́, позавчера́, ра́ньше, неда́вно, давно́ & Вчера́ я чита́л. \\
现在时 & 词干 ＋ 人称词尾 & сейча́с, сего́дня, у́тром, днём & Сейча́с я чита́ю. \\
将来时 & бу́ду ＋ 不定式 & за́втра, послеза́втра, ско́ро & За́втра я бу́ду чита́ть. \\
\end{tabular}
\tabbot
\end{minipage}

\textbf{这张表在说什么}：三行的动词部分完全不同，但时间副词是「配套的钥匙」——看到 вчера́ 就该用过去时，看到 за́втра 就该用将来时。

\textbf{逐行读法}：第二行的时间副词里 \textbf{сего́дня}（今天）既可与现在时也可与过去时搭配（今天做了），要看句子的意思；第三行的 \textbf{ско́ро}（很快）只配将来时。

\textbf{补：一句话的骨架（动过 p12）}。课件在这一页画了整整一张句子成分图，例子是 \textbf{Вчера́ у́тром там мой ста́рший брат до́лго писа́л дома́шнее зада́ние.}。这张图的价值是告诉你\textbf{俄语的语序不是「主谓宾」，而是「时间—地点—谁—怎么—做什么—做什么」}：

\begin{itemize}
\item 时间：вчера́ у́тром（昨天早上）
\item 地点：там（在那里）
\item 谁：мой ста́рший брат（我的哥哥）—— 物主 ＋ 形容词 ＋ 名词
\item 怎么：до́лго（很久）
\item 做什么：писа́л（写）
\item 做什么（宾语）：дома́шнее зада́ние（家庭作业）
\end{itemize}

所以动词永远在宾语前面、时间地点在句首——这条语序与 \S2.1 的「答句用 он／она́ 顶替名词」是配套的：两者都在说「俄语靠语序和代词，而不是靠汉语的「把」字结构」。

\subhdr{6.6　по-ру́сски 与副词：说、写、读、理解用什么语言}

动现 p26--p35 是这一讲的第二块内容，讲两件容易混的事——\textbf{语言名称做宾语}，和\textbf{по-ру́сски 这类副词做状语}：

\begin{minipage}{\linewidth}
\tabtop
\begin{tabular}{@{}>{\raggedright\arraybackslash}p{0.30\linewidth}>{\raggedright\arraybackslash}p{0.24\linewidth}>{\raggedright\arraybackslash}p{0.34\linewidth}@{}}
\thead{动词} & \thead{接什么} & \thead{课件例（动现 p26--p28）} \\
\end{tabular}
\tabmid
\begin{tabular}{@{}>{\raggedright\arraybackslash}p{0.30\linewidth}>{\raggedright\arraybackslash}p{0.24\linewidth}>{\raggedright\arraybackslash}p{0.34\linewidth}@{}}
знать（知道） & ＋ 语言名称（第四格） & Я зна́ю англи́йский язы́к. \\
учи́ть（学） & ＋ 语言名称（第四格） & Я учу́ ру́сский язы́к. \\
изуча́ть（研究、学习） & ＋ 语言名称（第四格） & Ты изуча́ешь ру́сский язы́к? \\
говори́ть／чита́ть／писа́ть／понима́ть & ＋ по-ру́сски 一类副词 & Студе́нты пи́шут по-ру́сски. \\
\end{tabular}
\tabbot
\end{minipage}

\textbf{这张表在说什么}：上面三个动词后面跟的是\textbf{名词}（язы́к，第四格），下面四个动词后面跟的是\textbf{副词}（по-ру́сски、по-англи́йски、по-кита́йски）。所以不能说 \textbf{*говорю́ ру́сский язы́к}，也不能说 \textbf{*зна́ю по-ру́сски}——前者把副词换成了名词，后者把名词换成了副词。

\textbf{逐行读法}：第一行 знать 是「会（一门语言）」，第二行 учи́ть 是「学（背单词、练语法）」，第三行 изуча́ть 更正式，多用于学校课程；第四行四个动词都是「用某种语言做某事」，介词 по- 后面跟语言名的形容词形式加 -ски。

\aside{这条区分与 \S5.4 的 что／потому́ что 是同一类问题：\textbf{动词决定后面接什么}。здесь 记住的是「接名词还是接副词」这一层，比记四个动词更省力。}

\textbf{副词的两组}（动现 p27、p32）：

\begin{itemize}
\item \textbf{评价类}：отли́чно（极好）、хорошо́（好）、непло́хо（不错）、немно́го（一点）、пло́хо（差）。例：\textbf{Я отли́чно зна́ю кита́йский язы́к.} ／ \textbf{Она́ пло́хо зна́ет ру́сский язы́к.}
\item \textbf{方式类（成对的反义）}：ти́хо（轻声）— гро́мко（大声）、ме́дленно（慢）— бы́стро（快）、пра́вильно（对）— непра́вильно（错）、внима́тельно（仔细）— невнима́тельно（马虎）。例：\textbf{Преподава́тель гро́мко спра́шивает.} ／ \textbf{Студе́нты внима́тельно чита́ют текст и пра́вильно отвеча́ют.}
\end{itemize}

\textbf{语序}（动现 p29）：\textbf{Кто ＋ 副词 ＋ 动词 ＋ 副词}——\textbf{Я отли́чно говорю́ по-ру́сски.} 第一个副词说「多好」，第二个副词说「用哪种语言」。对话里最常用的两问两答：\textbf{Ты уже́ хорошо́ понима́ешь по-ру́сски? — Да, я уже́ хорошо́ понима́ю по-ру́сски.} ／ \textbf{— Нет, я ещё пло́хо понима́ю по-ру́сски.} ／ \textbf{Как вы говори́те по-англи́йски? — Мы непло́хо говори́м по-англи́йски.}

\subhdr{6.7　练习：六个小题}

\begin{enumerate}[label=\textbf{第 \arabic* 题}]
\item （动现 p5）写出现在时六个人称：понима́ть。
\item （动现 p5）修正：*Я писа́ю письмо́ → \blank。
\item （动过 p4）填过去时：（Вчера́ моя́ подру́га）\blank кни́гу。（чита́ть）
\item （动将 p8）变将来时：Сейча́с я чита́ю. → За́втра я \blank чита́ть.
\item （动现 p29）组句：я, отли́чно, говори́ть, по-ру́сски。
\item （动现 p34）填反义副词：Студе́нт отвеча́ет гро́мко, а студе́нтка \blank。
\end{enumerate}

\textbf{第 1 题}　关键一步：понима́ть 是 -ать（第一组），词干以元音结尾，六个词尾照抄 чита́ть。答案：понима́ю ｜ понима́ешь ｜ понима́ет ｜ понима́ем ｜ понима́ете ｜ понима́ют。

\textbf{第 2 题}　关键一步：писа́ть 词干末尾是 ш，第一人称写 -у。答案：Я пишу́ письмо́。

\textbf{第 3 题}　关键一步：主语 подру́га 是阴性 → -ла。答案：чита́ла。

\textbf{第 4 题}　关键一步：将来时＝бу́ду ＋ 不定式，动词不变位。答案：бу́ду。

\textbf{第 5 题}　关键一步：语序是「谁 ＋ 多好 ＋ 动词 ＋ 用哪种语言」。答案：Я отли́чно говорю́ по-ру́сски。

\textbf{第 6 题}　关键一步：在 6.6 的方式类副词里找 гро́мко 的反义面。答案：ти́хо。
